% Part: sets-functions-relations
% Chapter: relations
% Section: special-properties

\documentclass[../../../include/open-logic-section]{subfiles}

\begin{document}

\olfileid{sfr}{rel}{prp}
\olsection{Sipat Khusus Relasi}

\begin{intro}
Sawetara jinis relasi kerep banget ditemoni nganti diwenehi
jeneng khusus. Tuladhane, $\le$ lan~$\subseteq$
nggayutake anggotane domain dhewe-dhewe
(umpamane, $\Nat$ kanggo~$\le$ lan
$\Pow{A}$ kanggo~$\subseteq$) kanthi cara kang padha.
Kanggo mangerteni kanthi persis kepriye relasi kasebut
padha lan bedane, kita nggolongake relasi
miturut sawetara sipat khusus kang bisa diduweni.
Pranyata sawetara sipat khusus iki (utawa gabungane)
wigati banget: tatanan lan relasi ekuivalensi.
\end{intro}

\begin{defn}[Refleksivitas]
Relasi $R \subseteq A^2$ iku \emph{refleksif} yen lan mung yen,
kanggo saben $x \in A$, $Rxx$.
\end{defn}

\begin{defn}[Transitivitas]
Relasi $R \subseteq A^2$ iku \emph{transitif} yen lan mung yen,
manawa $Rxy$ lan $Ryz$ lumaku, $Rxz$ uga lumaku.
\end{defn}

\begin{defn}[Simetri]
Relasi~$R \subseteq A^2$ iku \emph{simetris} yen lan mung yen,
manawa $Rxy$ lumaku, $Ryx$ uga lumaku.
\end{defn}

\begin{defn}[Antisimetri]
Relasi~$R \subseteq A^2$ iku \emph{antisimetris} yen lan mung yen,
manawa $Rxy$ lan $Ryx$ lumaku bebarengan, mula $x=y$
(utawa, kanthi tembung liya: yen $x\neq y$,
mula $\lnot Rxy$ utawa $\lnot Ryx$).
\end{defn}

\begin{explain}
Ing relasi simetris, $Rxy$ lan $Ryx$ tansah lumaku
bebarengan, utawa loro-lorone ora lumaku.
Ing relasi antisimetris, siji-sijine cara supaya $Rxy$
lan $Ryx$ lumaku bebarengan yaiku yen $x = y$.
Gatekna yen syarat iki ora \emph{mbutuhake} $Rxy$ lan $Ryx$
lumaku nalika $x = y$; syarat iki mung ora nglarang bab kasebut.
Mula relasi antisimetris bisa refleksif, nanging ora
saben relasi antisimetris iku refleksif.
Uga gatekna yen antisimetris lan mung ora simetris iku
syarat kang beda. Nyatane, relasi bisa simetris lan
antisimetris bebarengan (tuladhane, relasi identitas).
\end{explain}

\begin{defn}[Koneksitas]
Relasi $R \subseteq A^2$ iku \emph{koneks} yen kanggo kabeh
$x,y\in A$, yen $x \neq y$, mula $Rxy$ utawa~$Ryx$.
\end{defn}

\begin{prob}
Wenehana tuladha relasi kang (a) refleksif lan simetris
nanging ora transitif, (b) refleksif lan antisimetris,
(c) antisimetris, transitif, nanging ora refleksif, lan
(d) refleksif, simetris, lan transitif.
Aja nganggo relasi ing wilangan utawa himpunan.
\end{prob}

\begin{defn}[Irrefleksivitas]
Relasi $R \subseteq A^2$ diarani \emph{irrefleksif}
yen, kanggo kabeh $x \in A$, $Rxx$ ora lumaku.
\end{defn}

\begin{defn}[Asimetri]
Relasi $R \subseteq A^2$ diarani \emph{asimetris}
yen, kanggo saben pasangan $x,y\in A$, $Rxy$ lan~$Ryx$
ora tau lumaku bebarengan.
\end{defn}

Gatekna yen $A \neq \emptyset$, mula ora ana relasi
irrefleksif ing~$A$ kang refleksif, lan saben relasi
asimetris ing~$A$ uga antisimetris.
Nanging, ana $R \subseteq A^2$ kang ora refleksif lan
uga ora irrefleksif, lan ana relasi antisimetris
kang ora asimetris.

\end{document}
